% Chapter 17: Models for robust inference
% Bayesian Data Analysis - Gelman et al.

\chapter{稳健推断模型}

\section{本章导论}

在前面的章节中，我们主要依赖正态分布、二项分布和泊松分布，以及这些分布的层次组合来建模数据和参数。然而，使用有限的分布类别会导致有限的、甚至不恰当的推断。许多问题超出了便捷模型的适用范围，模型的选择应该拟合底层科学和数据，而不是仅仅为了分析或计算上的方便。正如第 5 章所说明的，创建现实模型最有用的方法通常是层次化地工作，将简单的单变量模型组合起来。

如果我们为了方便而使用过于简化的模型，就必须回答一个关键问题：后验推断在何种程度上依赖于极端数据点和不可评估的模型假设？我们在第 6 章已经讨论过这个问题的后半部分，这本质上是敏感性分析的主题；现在我们用更先进的计算方法来更深入地探讨这个话题。

\textbf{历史脉络}。贝叶斯稳健推断的现代发展可以追溯到 Box 和 Tiao 的开创性工作（Box and Tiao, 1968, 1973），他们系统地研究了正态模型中异常值的贝叶斯处理方法。这一传统后来被 Berger 和同事们进一步发展，探索了贝叶斯稳健性的理论方面，包括最大稳健性先验族等。

\section{稳健性的各个方面}

\subsection{对异常值的稳健性}

基于正态分布的模型对异常值是出了名的``非稳健''——在某种意义上，一个异常数据点可以强烈影响模型中所有参数的推断，即使那些参数与该异常观测没有实质性关联。

例如，在第 5.5 节的学业测试例子中，我们对八个处理效应的估计是通过将各学校的样本均值向总体均值收缩（即向真实效应来自共同正态分布的先验信息收缩）而得到的，每所学校 $j$ 的收缩比例仅由其抽样误差 $\sigma_j$ 和学校间效应变异 $\tau$ 决定。

假设研究中第八所学校的观测值 $y_8$ （表 5.2 中报告的）本应是 100 而不是 12，那么八个观测值变为 28、8、−3、7、−1、1、18 和 100，标准误与表 5.2 中报告的相同。如果我们对这一数据集应用层次正态模型，后验分布会告诉我们 $\tau$ 有一个很高的值，因此每个估计值 $\hat{\theta}_j$ 将基本上等于其观测效应 $y_j$。

但这在实践中合理吗？毕竟，给定这些假设观测值，第八所学校似乎有一个极其有效的 SAT 辅导项目，或者 100 可能只是数据记录错误的结果。在这两种情况下，让单个观测值 $y_8$ 对我们估计 $\theta_1, \ldots, \theta_7$ 的方式产生如此强烈的影响，似乎是不合适的。

在贝叶斯框架中，我们可以通过将 $\theta_j$ 的正态总体模型替换为更长尾的分布族来减少异常第八观测值的影响。这里所说的长尾分布是指在远离中心的位置有相对较高的概率含量，其中``远离''的尺度是相对于包含分布 50\% 概率区域的直径来确定的。长尾分布的例子包括 $t$ 分布族（其中最极端的情况是柯西分布或 $t_1$），以及（有限）混合模型——后者通常使用简单分布（如正态分布）来拟合大部分值，但允许来自替代分布的观测或参数值有离散的将概率，该替代分布可以有不同的中心，通常有更大的散布。

在这个假设的 SAT 辅导例子修改中，使用长尾分布对 $\theta_j$ 进行分析会导致观测值 100 被解释为来自长尾分布的极端抽取，而不是正态分布效应具有高方差的证据。结果是，分析会将第八个观测值在一定程度上向其他观测值收缩，但相对于其与总体均值的距离而言，收缩程度远不如前七个相互之间的收缩程度那么大。

正如我们的假设例子所示，我们不必为了处理异常值而放弃贝叶斯原理。例如，像柯西分布甚至双组分混合模型这样的长尾模型仍然是 $(\theta_1, \ldots, \theta_8)$ 的可交换先验分布——这在没有什么先验信息区分八所学校时是恰当的。可交换先验模型的选择影响 $\theta_j$ 估计值的收缩方式，我们因此可以在不完全改变分析方式的情况下减少异常观测值的影响。（这不应该取代对数据的仔细检查和对异常值中可能存在的记录错误的查找。）有时会区分搜索异常值的方法（可能将其从分析中移除）和使推断不受异常值影响的稳健方法。在贝叶斯框架中，这两种方法不应该被区分对待。例如，使用混合模型（第 22 章中的有限混合模型或标准模型的过度分散版本）不仅能将极端观测分类为来自高方差混合成分（而不是简单地将其视为令人惊讶的``异常值''），还意味着这些点对总体均值和中位数等估计量的推断影响较小。

\subsection{敏感性分析}

除了补偿异常值外，稳健模型还可用于评估后验推断对模型假设的敏感性。例如，可以使用 $t$ 分布替代正态分布来进行敏感性分析，以评估对正态假设的敏感性，方法是将自由度从大变小。如第 6 章所讨论的，敏感性分析的基本思想是尝试各种不同的分布（对于似然和先验模型），并观察感兴趣估计量和预测量的后验推断如何变化。一旦已经从后验分布中抽取了样本，使用重要性重采样可以相当直接地从替代模型中抽取样本，精度足以检测模型间推断的主要差异。如果感兴趣估计量的后验分布对模型假设高度敏感，则可能需要迭代模拟方法来获得更精确的计算。

从某种意义上说，第 5.5 节中 SAT 辅导实验的大部分分析，特别是图 5.6 和 5.7，是一种敏感性分析，其中参数 $\tau$ 被允许从 0 变化到 $\infty$。如第 5.5 节所讨论的，观测数据实际上与所有相等效应模型（即 $\tau = 0$）一致，但该模型没有任何实质性意义，所以我们拟合允许 $\tau$ 为任意正值的模型。结果总结在 $\tau$ 的边际后验分布中（如图 5.5 所示），描述了数据支持的 $\tau$ 值范围。

\section{标准概率模型的过度分散版本}

有时使用标准模型之一——二项、正态、泊松、指数——看起来很自然，但数据的分散程度过大。例如，正态分布不应该用于拟合大样本中 10\% 的点距离中位数超过四分位距 1.5 倍的数据。在上一节的假设例子中，我们建议 $\theta_j$ 的先验或总体模型应该有比正态分布更长的尾巴。实际上，对于每个标准模型，都有一个自然的扩展，其中添加一个参数来允许过度分散。每个扩展模型都有作为混合分布的解释。

\textbf{关于过度分散的一个关键性质}。所有这些分布的一个特征是：基于混合表示，它们永远不可能欠分散。\footnote{推导见附录 \cref{sec:overdispersion-derivation}} 具体而言，对于正态-逆-$\chi^2$ 混合（产生 $t$ 分布）、Gamma-Poisson 混合（产生负二项分布）和 Beta-Binomial 混合，每种情况下的边际方差都等于或大于对应的原始分布的方差。这是混合结构的直接结果：条件均值的条件方差永远非负，因此边际方差 = 条件均值的无条件方差 + 条件方差的条件期望，前者贡献了额外的分散。相比之下，如果数据被认为相对于标准分布欠分散，则应使用不同的模型（如零膨胀模型或过分散的离散分布的反面）。

\subsection{$t$ 分布替代正态分布}

$t$ 分布比正态分布有更长的尾巴，可用于适应（1）数据分布中偶尔出现的异常观测值，或（2）先验分布或层次模型中偶尔出现的极端参数。$t$ 分布族——$t_\nu(\mu, \sigma^2)$——由三个参数表征：中心 $\mu$、尺度 $\sigma$ 和决定分布形状的``自由度''参数 $\nu$。$t$ 密度是对称的，$\nu$ 必须在 $(0, \infty)$ 范围内。当 $\nu = 1$ 时，$t$ 等价于柯西分布，其尾巴如此之长以至于均值和方差都不存在；当 $\nu \to \infty$ 时，$t$ 趋近于正态分布。如果 $t$ 分布是试图准确拟合长尾分布的概率模型的一部分，基于相当大量的数据，那么通常将自由度作为未知参数是恰当的。在选择 $t$ 作为正态分布的稳健替代方案的应用中，自由度可以固定在一个小值上以允许异常值，但不能小于先验理解所规定的值。例如，自由度为 1 或 2 的 $t$ 分布方差无限，在远尾处通常不现实。

\textbf{混合解释}。回想第 3.2 节和第 12.1 节，$t_\nu(\mu, \sigma^2)$ 分布可以解释为具有共同均值和方差服从逆-$\chi^2$ 分布的正态分布的混合。\footnote{推导见附录 \cref{sec:t-distribution-mixture-derivation}} 例如，模型 $y_i \sim t_\nu(\mu, \sigma^2)$ 等价于
\[
y_i \mid V_i \sim \mathrm{N}(\mu, V_i), \quad V_i \sim \mathrm{Inv-}\chi^2(\nu, \sigma^2).
\label{eq:t-distribution-mixture}
\]
这是我们在第 12.1 节中作为辅助变量和参数扩展的计算方法说明而引入的。从统计上讲，具有高方差的观测值可以被视为分布中的异常值。类似地，当建模可交换参数 $\theta_j$ 时也有类似的解释。

\textbf{不同 $\nu$ 值的直观理解}。为了建立直觉，考虑 $\nu$ 取不同值时 $t$ 分布的形状变化：$\nu = 1$（柯西）有极长尾；$\nu = 2$ 方差仍无限但尾巴收敛更快；$\nu = 4$ 是常用的稳健选择，接近 logistic 形状；$\nu = 10$ 已非常接近正态但仍有稍厚尾；$\nu \to \infty$ 精确等于正态分布。这种连续过渡使得敏感性分析可以系统地探测尾部行为的影响。

\subsection{负二项分布替代泊松分布}

在将泊松模型应用于数据时，一个常见困难是泊松模型要求方差等于均值；而在实践中，计数数据的分布通常过度分散，方差大于均值。我们已经在广义线性模型的背景下讨论过过度分散问题（第 16.1 节），第 16.4 节给出了层次正态模型用于过度分散泊松回归的例子。

另一种对过度分散计数数据进行建模的方法是使用负二项分布，这是一个允许均值和方差分别拟合的二参数族。\footnote{推导见附录 \cref{sec:negative-binomial-variance-derivation}} 服从 $\mathrm{Neg-bin}(\alpha, \beta)$ 分布的数据 $y_1, \ldots, y_n$ 可以被认为是具有均值 $\lambda_1, \ldots, \lambda_n$ 的泊松观测值，这些均值服从 $\mathrm{Gamma}(\alpha, \beta)$ 分布。负二项分布的方差是
\[
\mathrm{var}(y) = \beta + \frac{\beta^2}{\alpha}.
\label{eq:negative-binomial-variance}
\]
总是大于均值 $\beta$；这与泊松分布的方差总是等于其均值形成对比。当 $\alpha \to \infty$ 时，底层 gamma 分布趋近于一个确定性的尖峰，负二项分布退化为泊松分布。

\subsection{贝塔-二项分布替代二项分布}

类似地，二项离散数据模型的实际局限在于只有一个自由参数，这意味着方差由均值决定。标准的稳健替代方案是贝塔-二项分布（beta-binomial），顾名思义，它是二项分布的贝塔混合。贝塔-二项分布用于建模教育测试数据等，其中``成功''是正确答案，个人正确回答的概率差异很大。这里，数据 $y_i$——每个人 $i = 1, \ldots, n$ 的正确答案数——用 $\mathrm{Beta\text{-}bin}(m, \alpha, \beta)$ 分布建模，被认为是具有共同试验次数 $m$ 和不等概率 $\pi_1, \ldots, \pi_n$ 的二项观测值，这些概率服从 $\mathrm{Beta}(\alpha, \beta)$ 分布。贝塔-二项分布的方差大于具有相同均值的二项分布的方差，额外因子为 $\frac{\alpha+\beta+m}{\alpha+\beta+1}$。

\subsection{logistic 和 probit 回归的 $t$ 分布替代}

logistic 和 probit 回归可以说是不稳健的，因为线性预测器 $X\beta$ 的绝对值较大时，逆 logit 或 probit 变换给出的概率接近 0 或 1。可以通过允许对 $X\beta$ 的大值偶尔出现错误预测来使这些模型更加稳健。这种稳健性不是针对数据 $y$（在二元回归中等于 0 或 1）来定义的，而是针对预测变量 $X$ 定义的。一个更稳健的模型允许离散回归模型拟合大部分数据，同时偶尔产生孤立误差。

一个稳健模型——robit 回归——可以使用离散数据回归模型的潜变量公式（第 408 页）来实现，将潜在连续数据 $u$ 的 logistic 或正态分布替换为模型 $u_i \sim t_\nu((X\beta)_i, 1)$。在现实设置中，从数据中估计 $\nu$ 是不切实际的——因为潜数据 $u_i$ 从未被直接观测到，因此基本上不可能形成关于其连续底层分布形状的推断——所以将其设定为一个低值以确保稳健性。设定 $\nu = 4$ 得到的分布接近 logistic，当 $\nu \to \infty$ 时，模型趋近于 probit。

\subsection{为什么还要使用非稳健模型？}

$t$ 族包括正态分布作为特例，那么为什么我们还要使用正态分布，或者二项、泊松或其他标准模型呢？首先，每个标准模型都有一个逻辑状态，使其对许多应用问题是合理的。二项和多项分布适用于独立同分布结果的离散计数，具有固定的计数总数。泊松和指数分布拟合泊松过程的事件数和等待时间，这是按时间索引的独立离散事件的自然模型。最后，中心极限定理告诉我们，正态分布是作为大量独立分量之和形成的数据的适当模型。

即使它们不是自然由问题结构暗示的，标准模型在计算上也是方便的，因为共轭先验分布通常允许直接计算后验均值和方差以及容易进行模拟。这就是为什么对学业测试例子中的 $\theta_j$ 拟合正态总体模型很容易，以及为什么通常对全正值数据的对数或限制在 0 到 1 之间数据的 logit 拟合正态模型。这里需要做一个重要区分：当标准模型在结构上合理时，使用它们是正当的；但当模型是为了计算方便而任意选择时，敏感性分析就成为必要的步骤——而不是表明稳健模型不重要。在后一种情况下，稳健模型的使用与其说是``替代''，不如说是对假设的``压力测试''。当模型检查（第六章）显示标准模型拟合良好时，我们可以更有信心地使用它；但即使如此，探索稳健替代方案也能帮助我们理解推断对尾部假设的敏感性。

\section{后验推断与计算}

一如既往，我们可以使用第三部分描述的方法从后验分布（在敏感性分析情况下是一个或多个分布）中进行抽样。在本节中，我们简要描述在混合公式下稳健模型的吉布斯抽样使用。该方法在第 17.4 节的层次正态-$t$ 模型中进行了说明。然而，在扩展模型时，我们有可能使用不那么耗时的近似方法：我们可以用原始后验分布的抽取作为新模型模拟的起点。在本节中，我们还描述了可用于稳健模型和敏感性分析的两种技术：用于计算边际后验密度的重要性加权，以及用于近似稳健分析的重要性重采样。

\subsection{稳健模型作为简单模型扩展的符号}

我们使用 $p_0(\theta \mid y)$ 表示原始模型的后验分布，假设已经拟合到数据，使用 $\varphi$ 表示表征用于稳健性或敏感性分析的扩展模型的超参数。我们的目标是从
\[
p(\theta \mid \varphi, y) \propto p(\theta \mid \varphi) \, p(y \mid \theta, \varphi).
\label{eq:robust-extension-posterior}
\]
中抽样，使用 $\varphi$ 的预指定值（如稳健 $t$ 模型的 $\nu = 4$），或者对于一系列 $\varphi$ 值。在后一种情况下，我们还希望计算敏感性分析参数的边际后验分布 $p(\varphi \mid y)$。

稳健分布族可以通过 $p(\theta \mid \varphi)$ 或 $p(y \mid \theta, \varphi)$ 进入模型。例如，第 17.2 节专注于稳健数据分布，而我们第 17.4 节对 SAT 辅导实验的重新分析使用稳健分布建模模型参数。我们必须建立一个联合先验分布 $p(\theta, \varphi)$，这需要一些注意，因为它捕捉了 $\theta$ 和 $\varphi$ 之间的先验依赖关系。

\subsection{使用混合公式的吉布斯抽样}

马尔可夫链模拟可用于从后验分布 $p(\theta \mid \varphi, y)$ 中抽取。这可以使用混合公式完成，方法是从 $\theta$ 和额外未观测尺度参数（$t$ 模型中的 $V_i$、负二项模型中的 $\lambda_i$、贝塔-二项模型中的 $\pi_i$）的联合后验分布中抽样。

一个简单的例子，考虑将 $t_\nu(\mu, \sigma^2)$ 分布拟合到数据 $y_1, \ldots, y_n$，其中 $\mu$ 和 $\sigma$ 未知。给定 $\nu$，我们已经在第 12.1 节讨论了如何将吉布斯抽样器编程为涉及 $\mu$、$\sigma^2$、$V_1, \ldots, V_n$ 的参数化形式。如果 $\nu$ 本身未知，吉布斯抽样器必须扩展到包括从 $\nu$ 的条件后验分布中抽样的步骤。没有简单的方法来进行这一步，但可以使用 Metropolis 步骤。另一个复杂之处在于这类模型通常有多峰后验密度，不同的模式对应于 $t$ 分布尾巴中不同观测值的模式，这意味着在初始阶段需要额外的工作来搜索模式并在模拟中在模式之间跳跃，例如使用模拟退火（见第 12.3 节）。

\section{八校例子中的稳健推断与敏感性分析}

考虑基于第 5.5 节表 5.2 中数据的 SAT 辅导效应的层次模型。鉴于八个原始实验的大样本量，假设数据模型 $y_j \mid \theta_j, \sigma_j^2 \sim \mathrm{N}(\theta_j, \sigma_j^2)$（方差 $\sigma_j^2$ 已知）应该没有太大问题。总体模型 $\theta_j \mid \mu, \tau^2 \sim \mathrm{N}(\mu, \tau^2)$ 更难证明，尽管第 6.5 节的模型检查表明它对于获得学校效应的后验区间是足够的。然而，一般来说，后验推断可能对假设模型高度敏感，即使模型对观测数据提供了良好的拟合。为了说明稳健推断和敏感性分析方法，我们探索一个替代模型族，将 $t$ 分布拟合到学校效应总体：
\[
\theta_j \mid \nu, \mu, \tau \sim t_\nu(\mu, \tau^2), \quad j = 1, \ldots, 8.
\label{eq:t-population-model}
\]

我们使用符号 $p(\theta, \mu, \tau \mid \nu, y) \propto p(\theta, \mu, \tau \mid \nu) \, p(y \mid \theta, \mu, \tau, \nu)$ 表示在 $t_\nu$ 模型下的后验分布，$p_0(\theta, \mu, \tau \mid y) \equiv p(\theta, \mu, \tau \mid \nu = \infty, y)$ 表示第 5.5 节评估的正态模型下的后验分布。

\subsection{基于 $t_4$ 总体分布的稳健推断}

正如本章开头所讨论的，人们可能会担心正态总体模型导致最极端估计的学校效应被过多地向总体均值收缩。也许某个学校的辅导项目与其他学校足够不同，其估计值不应该被过多地向平均值收缩。一个相关担忧是最大观测效应可能对总体方差 $\tau^2$ 的估计产生不当影响，进而影响其他效应的贝叶斯估计。从建模角度来看，各种 SAT 辅导项目差异很大，其效应的总体可能更好地用长尾分布来拟合。为了评估这些担忧的重要性，我们执行稳健分析，用 $t$ 模型（\cref{eq:t-population-model}）替换正态总体分布，设定 $\nu = 4$，其他模型保持不变。

\textbf{吉布斯抽样}。我们使用第 12.1 节描述的方法进行吉布斯抽样，设定 $\nu = 4$。\footnote{详细计算步骤见附录 \cref{sec:eight-schools-gibbs-derivation}} 结果推断（基于从后验分布中抽取的 2500 次模拟——五个链的最后一半，每个链长 1000）汇总在表 17.1 中。结果与第 5.3 节正态模型下获得的推断\textbf{在中心趋势（后验中位数、均值）上基本相同}，但稳健模型对极端学校（如学校 A）的收缩程度稍小——这正是稳健分析的预期行为：减少对异常值的过度收缩。

\textbf{使用重要性重采样的计算}。尽管我们已经完成了马尔可夫链模拟，我们简要讨论如何应用重要性重采样来近似 $\nu = 4$ 的后验分布。首先，我们从 $p_0(\theta, \mu, \tau \mid y)$（正态模型下的后验分布）中抽取 5000 个 $(\theta, \mu, \tau)$ 的抽取值。接下来，我们计算每个抽取值的重要性比率：\footnote{推导见附录 \cref{sec:importance-ratio-derivation}}
\[
\frac{p(\theta, \mu, \tau \mid \nu, y)}{p_0(\theta, \mu, \tau \mid y)} = \frac{p(\mu, \tau \mid \nu) \, p(\theta \mid \mu, \tau, \nu) \, p(y \mid \theta, \mu, \tau, \nu)}{p_0(\mu, \tau) \, p_0(\theta \mid \mu, \tau) \, p_0(y \mid \theta, \mu, \tau)} \propto \prod_{j=1}^8 \frac{t_\nu(\theta_j \mid \mu, \tau^2)}{\mathrm{N}(\theta_j \mid \mu, \tau^2)}.
\label{eq:importance-ratio}
\]

然后我们使用重要性重采样，从 5000 个抽取值中无放回地抽取 500 个 $(\theta, \mu, \tau)$ 的子样本。这种近似对于评估稳健性可能已足够，但重要性比率对数的长尾分布（未显示）确实表明使用重要性重采样获得准确推断存在严重问题。

\subsection{基于不同 $\nu$ 值的 $t_\nu$ 分布的敏感性分析}

与稳健性略有不同的担忧是对正态总体分布先验假设的后验推断敏感性。为了研究敏感性，我们现在拟合一系列 $t$ 分布，自由度分别为 1、2、3、5、10 和 30。我们已经拟合了无限自由度（正态模型）和 4 个自由度（上述稳健模型）。

对于每个 $\nu$ 值，我们执行马尔可夫链模拟以从 $p(\theta, \mu, \tau \mid \nu, y)$ 中获取抽取值。我们通过每个学校效应 $\theta_j$ 的后验均值和标准差来总结结果。\cref{fig:nu-sensitivity-plot} 显示了作为 $\nu$ 函数的这些结果。在 $\nu_1$ 而非 $\nu$ 参数化方面有优势，因为包括正态分布在 $\nu_1 = 0$ 处，并将从正态到柯西分布的整个范围包含在有限区间 $[0, 1]$ 中。图中有一些变化，但没有显示出对超参数的系统性敏感性。

\subsection{将 $\nu$ 作为未知参数处理}

最后，我们将敏感性分析参数 $\nu$ 视为未知量，并在后验分布中对其进行平均。一般来说，这是关键的一步，因为我们通常只关心被数据支持的模型的敏感性。在这个特定例子中，推断对 $\nu$ 不敏感，因此计算边际后验分布是不必要的；我们在这里将其作为通用方法的说明。

\textbf{先验分布}。在计算 $\nu$ 的后验分布之前，我们必须为其指定先验分布。我们尝试在 $\nu_1$ 上使用均匀密度，范围为 $[0, 1]$（即从正态到柯西分布）。这个先验分布偏向长尾模型，一半的先验概率落在 $t_1$（柯西）和 $t_2$ 分布之间。

\textbf{后验推断}。为了将 $\nu$ 作为未知参数处理，我们修改稳健分析中使用的吉布斯抽样模拟，包括一个 Metropolis 步骤用于从 $\nu_1$ 的条件分布中抽样。\cref{fig:nu-posterior-histogram} 显示了 $\nu_1$ 模拟的直方图。

\textbf{讨论}。稳健性和对建模假设的敏感性取决于所研究的估计量。在 SAT 辅导例子中，八个学校效应的后验中位数、50\% 和 95\% 区间对正态总体分布假设不敏感（至少与 $t$ 族相比）。相比之下，99.9\% 区间可能强烈依赖于分布的尾部，并对 $t$ 分布中的自由度敏感——幸运的是，这些极端尾部在这个例子中不太可能有实质性意义。这一发现的意义在于：我们不需要担心正态假设会严重影响我们对这个特定数据集的中心推断；但如果我们对极尾部（如某效应超过 50 分的概率）感兴趣，那么选择哪个 $t$ 分布就变得关键了。

\section{使用 $t$ 分布误差的稳健回归}

与基于正态分布的其他模型一样，第 14 章正态线性回归模型下的推断对异常值或极端值敏感。稳健回归分析通过考虑回归误差的正态分布的稳健替代方案（如自由度数较少的 $t$ 分布）来获得。稳健误差分布（如方差较少的 $t$ 分布）将远离回归线的观测视为高方差观测，得到的结果类似于对异常值降权的结果。

\subsection{迭代加权线性回归和 EM 算法}

为了说明稳健回归计算，我们考虑将 $t_\nu$ 回归模型作为正态线性回归模型的替代方案，固定自由度参数。给定解释变量向量 $X_i$，个体响应变量 $y_i$ 的条件分布是 $p(y_i \mid X_i \beta, \sigma^2) = t_\nu(y_i \mid X_i \beta, \sigma^2)$。$t_\nu$ 分布可以表示为混合形式。$t$ 模型下 $p(\beta, \sigma^2 \mid \nu, y)$ 后验分布的模式可以直接使用牛顿法（第 13.1 节）或其他模式寻找技术获得。\footnote{详细推导见附录 \cref{sec:robust-regression-em-derivation}} 或者，我们可以利用 $t$ 模型的混合形式，使用 EM 算法，将方差 $V_i$ 视为``缺失数据''（即要平均的参数）来处理。

E 步计算在当前参数估计值给定的情况下，充分统计量的期望值。\footnote{详见附录 \cref{sec:robust-regression-em-derivation}} M 步是具有对角权重矩阵 $W$ 的加权线性回归，其中权重矩阵对角线包含 $1/V_i$ 的条件期望。EM 算法的迭代等价于迭代加权最小二乘算法中执行的迭代。给定回归参数的初始估计，计算每个案例的权重，残差较大的案例权重较小。然后通过加权线性回归获得回归参数的改进估计。

当自由度参数 $\nu$ 被视为未知时，可以应用 ECME 算法，添加一个额外步骤来更新自由度。

\textbf{其他稳健模型}。迭代加权线性回归或等效的 EM 算法可用于获得许多稳健替代模型的后验模式。更改用于观测方差 $V_i$ 的概率模型会产生替代稳健模型。例如，可以使用两点分布来建模具有污染误差的回归。

\subsection{吉布斯抽样和 Metropolis 算法}

可以使用吉布斯抽样和 Metropolis 算法从稳健回归模型的后验分布中获取样本。使用 $t_\nu$ 分布的混合参数化，我们可以通过交替从 $p(\beta, \sigma^2 \mid V_1, \ldots, V_n, \nu, y)$（使用加权线性回归的通常后验分布）和从 $p(V_1, \ldots, V_n \mid \beta, \sigma^2, \nu, y)$（一组独立的缩放逆-$\chi^2$ 分布）中抽取来获得后验样本。

\section{本章小结}

本章我们学习了：

\begin{enumerate}
  \item \textbf{稳健性的概念}：后验推断对异常值和模型假设的敏感性，以及 Box-Tiao 的历史贡献
  \item \textbf{标准模型的稳健替代}：
  \begin{enumerate}
    \item $t$ 分布替代正态分布（可使用正态-逆-$\chi^2$ 混合表示）
    \item 负二项分布替代泊松分布（Gamma-Poisson 混合）
    \item 贝塔-二项分布替代二项分布
    \item robit 回归（$t$ 分布替代 logistic/probit）
  \end{enumerate}
  \item \textbf{后验计算方法}：
  \begin{enumerate}
    \item 使用混合公式的吉布斯抽样
    \item 重要性加权和重要性重采样
  \end{enumerate}
  \item \textbf{敏感性分析}：通过将 $\nu$ 作为自由度参数变化来评估对假设的敏感性
  \item \textbf{稳健回归}：使用 $t$ 分布误差的稳健线性回归
\end{enumerate}

\section{下节预告}

下一章我们将学习缺失数据模型，这是贝叶斯分析中一个重要的实际问题，涉及如何对观测不完整的数据进行建模和推断。

\endinput
